यदि $!A \ident (!B \lif !C)$, तो $\pi$ निम्न पर समाप्त होता है
\[
\AxiomC{}
\RightLabel{$\pi_1'$}
\Deduce$!B, \Gamma \fCenter \Delta, !C$
\RightLabel{\RightR{\lif}}
\UnaryInf$\Gamma \fCenter \Delta, !B \lif !C$
\LeftSubproofLabel{$\pi_1$}
\AxiomC{}
\RightLabel{$\pi_2'$}
\Deduce$!B \lif !C, \Pi \fCenter \Lambda, !B$
\AxiomC{}
\RightLabel{$\pi_2''$}
\Deduce$!C, \Pi \fCenter \Lambda$
\RightLabel{\LeftR{\lif}}
\BinaryInf$!B \lif !C, \Pi \fCenter \Lambda$
\RightSubproofLabel{$\pi_2$}
\RightLabel{\Cut}
\BinaryInf$\Gamma, \Pi \fCenter \Delta, \Lambda$
\DisplayProof
\]
\Cut पर समाप्त होने वाले !!{proof}
\[
\AxiomC{}
\RightLabel{$\pi_1'$}
\Deduce$!B, \Gamma \fCenter \Delta, !C$
\RightLabel{\RightR{\lif}}
\UnaryInf$\Gamma \fCenter \Delta, !B \lif !C$
\LeftSubproofLabel{$\pi_1$}
\AxiomC{}
\RightLabel{$\pi_2'$}
\Deduce$!B \lif !C, \Pi \fCenter \Lambda, !B$
\RightLabel{\Cut}
\BinaryInf$\Gamma, \Pi \fCenter \Delta, \Lambda, !B$
\DisplayProof
\]
की कट !!{height},~$\pi$ से कम है। आगमन परिकल्पना $\Gamma, \Pi \fCenter
\Delta, \Lambda, !B$ का !!a{proof}~$\pi_1''$ देती है। \Cut पर समाप्त होने
वाला !!{proof}
\[
\AxiomC{}
\RightLabel{$\pi_1'$}
\Deduce$B!, \Gamma \fCenter \Delta, !C$
\AxiomC{}
\RightLabel{$\pi_2''$}
\Deduce$!C, \Pi \fCenter \Lambda$
\RightLabel{\Cut}
\BinaryInf$!B, \Gamma, \Pi \fCenter \Delta, \Lambda$
\DisplayProof
\]
जिसकी कट कोटि और कट !!{height} दोनों~$\pi$ की कट कोटि और !!{height} से कम
हैं; अतः आगमन परिकल्पना $!B, \Gamma, \Pi, \Pi \fCenter \Delta, \Lambda,
\Lambda$ का \Log{G3c}-!!{proof}~$\pi_2'''$ देती है। अब \Cut पर समाप्त होने
वाले !!{proof}
\[
\AxiomC{}
\RightLabel{$\pi_1''$}
\Deduce$\Gamma, \Pi \fCenter \Delta, \Lambda, !B$
\AxiomC{}
\RightLabel{$\pi_2'''$}
\Deduce$!B, \Gamma, \Pi, \Pi \fCenter \Delta, \Lambda,
\Lambda$
\RightLabel{\Cut}
\BinaryInf$!B, \Gamma, \Pi \fCenter \Delta, \Lambda$
\DisplayProof
\]
(जिसकी कट कोटि $< \cutr{\pi}$ है) पर आगमन परिकल्पना लगाकर $\Gamma, \Gamma,
\Pi, \Pi, \Pi \fCenter \Delta, \Lambda, \Lambda, \Lambda$ का !!a{proof}
प्राप्त करें। \olref[seq][inv]{prop:G3c-cont-adm} लगाकर हमें $\Gamma, \Pi
\fCenter \Delta, \Lambda$ का !!{proof} $\pi'$ मिलता है।
